% R. Brown's translations of two letters to L. Breen, with his note and comments (cote 134-2): archive cote 134-2, site page 119 (PDF page 120).
% comments. Draft compared: none.
% AI-transcribed by Claude Opus 5.5 (claude-opus-5-5, adaptive thinking, effort high), Message Batches API, batch brown_full (msgbatch_01Md3i5yK1PxkpJn5GHfEExE), 2026-11-13; not yet checked by eye.
\scanpage{134-2/119}
\mspage{-3-}

in:

J.-L. Loday, 'Spaces with finitely many homotopy groups', J. Pure
Appl. Alg.  24(1982) 179 - 202.
The n-cat-groups (which may be described as n-fold categories in
the category of groups) are adequate for the description of
truncated homotopy types.  This is Loday's main result.

7. An analysis of the structure of non-abelian chain complex on a
length two normalised chain complex of a simplicial group is
given in:
D. Conduché, 'Modules croisés généralisés de longeur 2' (in
preparation).

8. By following through some of the equivalences in (4) above, we
can associate to an $\infty$-groupoid G \quad a simplicial T-complex NG, whose
underlying simplicial set could be called the \textit{nerve} of G.  A
direct description of NG is that $(\mathrm{NG})_{n}$ is the set of
$\infty$-groupoid maps  to G from $\sigma\underset{\approx}{\Delta}^{n}$, the homotopy $\infty$-groupoid of the
standard n-simplex $\Delta^{n}$ (i.e. $\sigma\underset{\approx}{\Delta}^{n}$ has associated crossed complex $\pi\underset{\approx}{\Delta}^{n}$,
the homotopy crossed complex of $\underset{\sim}{\Delta}^{n}$).
Now a definition of $\infty$-category is also given in (4)(d) above,
but the corresponding functor nerve: $\infty$-categories $\longrightarrow$ simplicial
sets is still  unknown.

9. An indication of the use of the ideas of (4) in non-abelian
situations is given in:

a) R. Brown and P.J. Higgins, 'Crossed complexes and non-%
abelian extensions', Int. Conf. on Category Theory, Gummersbach
(1981),  Springer L.N.M. (to appear).

b) R. Brown, 'Non-abelian cohomology and the homotopy classification
of maps', Proc. Conf, \textit{Méthodes d'algèbre homotopique en topologie},
Marseilles - Luminy (1982).  Astérixes (to appear)

\mspage{\hand{120}}
